\begin{lemma}
Fix \(n,m,k\in\mathbb N\), \(p\in\mathbb R\), a finite set
\(I\subseteq\mathrm{Fin}(n)\), parameters
\(\varepsilon,\varepsilon_1,\varepsilon_2\in\mathbb R\), and \(n_0\in\mathbb N\).
Assume
\[
\noderef{ParameterHierarchyEventualComparisons}
  (\varepsilon,\varepsilon_1,\varepsilon_2,n_0),
\qquad n_0<n,
\]
\[
m=\noderef{TwoBiteNaturalReducedVertexCount}(n),\qquad
p=\noderef{TwoBiteEdgeProbability}\!\left(\frac12,n\right),
\]
\[
k=\noderef{TwoBiteNaturalIndependenceScale}(1+\varepsilon,n),
\qquad |I|=k.
\]
Put
\[
M=\noderef{TwoBiteNaturalReducedVertexCount}(n),\qquad
p_0=\noderef{TwoBiteEdgeProbability}\!\left(\frac12,n\right),
\]
\[
L=\log n,\qquad
K=\noderef{TwoBiteNaturalIndependenceScale}(1+\varepsilon,n),
\qquad
C=\noderef{RealChooseTwo}(K).
\]
Here \(K\) is viewed as a real number in the formula for \(C\).  For a
two-bite sample \(\omega\in\noderef{TwoBiteSample}(n,M,p_0)\), define
\[
P_R=\{\noderef{TwoBiteRedProjection}(\omega,u):u\in I\},\qquad
P_B=\{\noderef{TwoBiteBlueProjection}(\omega,u):u\in I\},
\]
and let \(P_I\) be the union of \(P_R\) in the red copy and \(P_B\) in the
blue copy of the base-vertex set.  Define
\[
S_I^{\mathrm{ext}}
=\{x:\noderef{TwoBiteSmallClass}(\omega,\varepsilon,I,x)\}\setminus P_I,
\]
and
\[
E_{\mathrm{ext}}(\omega)
=\bigcup_{x\in S_I^{\mathrm{ext}}}
  \noderef{TwoBiteFinalPairs}(\noderef{TwoBiteX}(\omega,I,x)).
\]
Let \(Z_I(\omega)\) be the number of pairs \(e=(u,v)\) in
\(E_{\mathrm{ext}}(\omega)\) satisfying
\[
\noderef{TwoBiteRedProjection}(\omega,u)
  \ne\noderef{TwoBiteRedProjection}(\omega,v),\qquad
\noderef{TwoBiteBlueProjection}(\omega,u)
  \ne\noderef{TwoBiteBlueProjection}(\omega,v).
\]
With \(\mu=C/(2L)\) and \(t=C/\sqrt L\),
\[
\noderef{TwoBiteEventProbability}(n,M,p_0)
 \bigl(\{\omega:Z_I(\omega)\ge \mu+t\}\bigr)
\le
\exp\!\left(
-\frac{C^2}{L\,M\,n^{8\varepsilon}}
\right).
\]
\end{lemma}
\begin{proof}
Fix the data in the statement and abbreviate \(M,p_0,L,K,C\) as above.  From
\(|I|=k\) and \(k=K\), we have \(|I|=K\).  The hierarchy assumptions also give
the range facts needed below: \(0\le p_0\le 1/2\), \(M\le n/L^2\), \(L>0\),
and \(n>0\).

Partition the event by the embedding
\(\pi=\noderef{TwoBiteEmbedding}(\omega)\).  For a fixed embedding \(\pi\),
apply \(\noderef{FixedSetProductModelIngredients}\) with this \(I\),
\(\varepsilon,\varepsilon_1,\varepsilon_2,n_0\), and \(\pi\).  It gives a
finite product space \(\alpha^{2M}\), a probability mass \(q\), and a product
variable \(Z_{\mathrm{prod}}\) such that
\[
q\ge0,\qquad \sum_{\alpha}q=1,
\]
\[
c=\frac12 n^{4\varepsilon}>0,\qquad t=\frac{C}{\sqrt L}\ge0,
\]
the one-coordinate change of \(Z_{\mathrm{prod}}\) is at most \(c\), and
\[
\noderef{TwoBiteEventProbability}(n,M,p_0)
 \bigl(\{\omega:\noderef{TwoBiteEmbedding}(\omega)=\pi
   \text{ and } Z_I(\omega)\ge C/(2L)+t\}\bigr)
\]
is at most
\[
\Pr_q\!\left(
Z_{\mathrm{prod}}\ge\mathbb E_qZ_{\mathrm{prod}}+t
\right)
\noderef{TwoBiteEventProbability}(n,M,p_0)
 \bigl(\{\omega:\noderef{TwoBiteEmbedding}(\omega)=\pi\}\bigr).
\]
The expectation comparison
\(\mathbb E_qZ_{\mathrm{prod}}\le C/(2L)\) is part of that helper.  It is what
allows the fixed threshold \(C/(2L)+t\) in the present event to be controlled
by the upper tail around the product expectation.

Apply \(\noderef{BoundedDifferencesProduct}\) to this product model, with
\[
r=2M,\qquad c=\frac12 n^{4\varepsilon},\qquad t=\frac{C}{\sqrt L}.
\]
The hypotheses just listed are exactly its mass, positivity, and
bounded-difference assumptions, so
\[
\Pr_q\!\left(
Z_{\mathrm{prod}}\ge\mathbb E_qZ_{\mathrm{prod}}+t
\right)
\le
\exp\!\left(-\frac{2t^2}{r c^2}\right).
\]
The exponent conversion
\(\noderef{FixedSetTailExponentInequality}\), for the same \(r,c,t\), gives
\[
-\frac{2t^2}{r c^2}
\le
-\frac{C^2}{L\,M\,n^{8\varepsilon}}.
\]
Since the exponential function is increasing, the preceding two displays give,
for every embedding \(\pi\),
\[
\begin{aligned}
&\noderef{TwoBiteEventProbability}(n,M,p_0)
 \bigl(\{\omega:\noderef{TwoBiteEmbedding}(\omega)=\pi
   \text{ and } Z_I(\omega)\ge C/(2L)+C/\sqrt L\}\bigr)\\
&\qquad\le
\exp\!\left(-\frac{C^2}{L\,M\,n^{8\varepsilon}}\right)
\noderef{TwoBiteEventProbability}(n,M,p_0)
 \bigl(\{\omega:\noderef{TwoBiteEmbedding}(\omega)=\pi\}\bigr).
\end{aligned}
\]

Finally apply \(\noderef{TwoBiteEventProbabilitySumOverEmbeddings}\).  The
embedding cells partition the finite sample space, and the bound on each cell
has the same nonnegative factor
\[
\exp\!\left(-\frac{C^2}{L\,M\,n^{8\varepsilon}}\right).
\]
Summing the cell estimates therefore bounds the unconditional probability of
\(Z_I\ge C/(2L)+C/\sqrt L\) by that factor, which is precisely the claimed
inequality.
\end{proof}
